\documentclass[10pt,letterpaper]{article}
\usepackage{cornell-notes}

\title{Clause 21.03.11: Constant Evaluation Context}
\author{}
\date{}
\setDocTitle{Clause 21.03.11: Constant Evaluation Context}
\setDocSubtitle{C++ 2024}
\setDocOwner{C++ 2024 Cornell Notes}
\setDocDate{\today}
\setCornellCollection{C++ 2024 Cornell Notes}
\setCornellUnitType{Clause}
\setCornellUnitNumber{21.03.11}
\setCornellUnitTitle{Constant Evaluation Context}
\hypersetup{pdftitle={Clause 21.03.11: Constant Evaluation Context},pdfsubject={Cornell notes},pdfauthor={}}

\begin{document}
\maketitle
\begin{CornellOverviewBox}{Overview}
\textbf{Classification:} Standard library - Normative\\[2pt]
\textbf{Learning objective:} Understand the rules, terminology, and semantic consequences associated with Constant evaluation context.
\end{CornellOverviewBox}\begin{CornellSummaryBox}{Summary}
{\sffamily\bfseries\color{Accent} Summary}\par\vspace{3pt}Section 21.3.11 focuses on Constant evaluation context. Apply the specified interface together with its mandates, effects, result conditions, exception behavior, and complexity constraints. When applying it, preserve the distinction between mandatory requirements, recommendations, permissions, and behavior deliberately left undefined, unspecified, or implementation-defined.
\end{CornellSummaryBox}

\end{document}
